% questions.tex -- 30 multiple-choice questions, Module 5 Topic 3 (gravity).
% Section A: Q1-6, Section B: Q7-18, Section C: Q19-30. Each section goes from easy to hard.
% Key (balanced A/B/C/D): 1C 2D 3A 4B 5D 6C 7A 8D 9D 10B 11C 12D 13A 14C 15A
%                         16B 17C 18D 19A 20B 21C 22D 23B 24C 25A 26C 27A 28B 29D 30B
% Every number: checks/m5t3-booklet/numbers.py
% \Needspace before each question keeps its stem, diagram and options together.
\unnumberedchapter[\faDumbbell]{Questions}

% =====================================================================
\section*{Section A: Gravitational force and field strength}

% ---- Q1 (Newton's third law)
\Needspace{12\baselineskip}
\begin{practiceq}{1}{}\label{bq:secA}
A satellite of mass \qty{1000}{kg} orbits Earth. Which statement about the gravitational forces
between Earth and the satellite is correct?
\choices[1]{Earth pulls on the satellite with a bigger force than the satellite pulls on Earth.}
  {The satellite does not pull on Earth, because its mass is too small.}
  {The two forces are equal in size, but Earth's acceleration is much smaller than the
   satellite's acceleration.}
  {The two forces are equal in size, so Earth and the satellite have equal accelerations.}
\end{practiceq}

% ---- Q2 (g from mass and radius ratios)
\Needspace{8\baselineskip}
\begin{practiceq}{1}{}
Planet P has 2 times the mass of Earth and half the radius of Earth. The gravitational field
strength at Earth's surface is $g$. What is the gravitational field strength at the surface of
planet P?
\choices{$\dfrac{g}{2}$}{$2g$}{$4g$}{$8g$}
\end{practiceq}

% ---- Q3 (altitude where g is one-ninth)
\Needspace{8\baselineskip}
\begin{practiceq}{1}{}
The gravitational field strength at Earth's surface is $g$. At what \emph{altitude} (height above
Earth's surface) is the gravitational field strength equal to $\dfrac{g}{9}$? The radius of Earth is
$r_\mathrm{E}$.
\choices{$2r_\mathrm{E}$}{$3r_\mathrm{E}$}{$8r_\mathrm{E}$}{$9r_\mathrm{E}$}
\end{practiceq}

% ---- Q4 (force of Mars on Phobos)
\Needspace{17\baselineskip}
\begin{practiceq}{1}{}
The table gives data for Mars and its moon Phobos.
\begin{center}
  \begin{tabular}{@{}L{70mm}R{32mm}@{}}
    \toprule
    Mass of Mars & \qty{6.42e23}{kg}\\
    Mass of Phobos & \qty{1.07e16}{kg}\\
    Mean orbital radius of Phobos & \qty{9.38e6}{m}\\
    Orbital period of Phobos & \qty{2.75e4}{s}\\
    \bottomrule
  \end{tabular}
\end{center}
What is the size of the gravitational force of Mars on Phobos?
\choices{\qty{0.49}{N}}{\qty{5.2e15}{N}}{\qty{5.2e21}{N}}{\qty{4.9e22}{N}}
\end{practiceq}

% ---- Q5 (pendulum on Earth and the Moon: T^2 against l)
\Needspace{28\baselineskip}
\begin{practiceq}{1}{}
A simple pendulum can be used to measure the gravitational field strength $g$, using
\[ T = 2\pi\sqrt{\dfrac{l}{g}} \]
where $T$ is the period and $l$ is the length of the pendulum. On the Moon, $g$ is about
one-sixth of its value on Earth.

A student plots $T^2$ (vertical axis) against $l$ (horizontal axis) for pendulums on Earth and on
the Moon. Which graph is correct?
\begin{center}
  Key:\quad
  \tikz[baseline=-0.6ex]\draw[hscseriesA, line width=1.3pt] (0,0) -- (0.9,0);\ Earth data
  \qquad
  \tikz[baseline=-0.6ex]\draw[hscseriesB, line width=1.3pt, dashed] (0,0) -- (0.9,0);\ Moon data
\end{center}
\choices[2]{%
  \begin{tikzpicture}[baseline=(current bounding box.north)]
    \draw[axisline] (0,0) -- (3.0,0) node[below left, font=\small] {$l$};
    \draw[axisline] (0,0) -- (0,2.2) node[above, font=\small] {$T^2$};
    \draw[hscseriesA, line width=1.3pt] (0,0) -- (2.6,1.85);
    \draw[hscseriesB, line width=1.3pt, dashed] (0,0) -- (2.6,0.35);
  \end{tikzpicture}}
  {%
  \begin{tikzpicture}[baseline=(current bounding box.north)]
    \draw[axisline] (0,0) -- (3.0,0) node[below left, font=\small] {$l$};
    \draw[axisline] (0,0) -- (0,2.2) node[above, font=\small] {$T^2$};
    \draw[hscseriesA, line width=1.3pt, domain=0:2.6, samples=60, smooth] plot (\x, {0.22*sqrt(\x)});
    \draw[hscseriesB, line width=1.3pt, dashed, domain=0:2.6, samples=60, smooth] plot (\x, {1.15*sqrt(\x)});
  \end{tikzpicture}}
  {%
  \begin{tikzpicture}[baseline=(current bounding box.north)]
    \draw[axisline] (0,0) -- (3.0,0) node[below left, font=\small] {$l$};
    \draw[axisline] (0,0) -- (0,2.2) node[above, font=\small] {$T^2$};
    \draw[hscseriesA, line width=1.3pt, domain=0:2.6, samples=60, smooth] plot (\x, {1.15*sqrt(\x)});
    \draw[hscseriesB, line width=1.3pt, dashed, domain=0:2.6, samples=60, smooth] plot (\x, {0.22*sqrt(\x)});
  \end{tikzpicture}}
  {%
  \begin{tikzpicture}[baseline=(current bounding box.north)]
    \draw[axisline] (0,0) -- (3.0,0) node[below left, font=\small] {$l$};
    \draw[axisline] (0,0) -- (0,2.2) node[above, font=\small] {$T^2$};
    \draw[hscseriesA, line width=1.3pt] (0,0) -- (2.6,0.35);
    \draw[hscseriesB, line width=1.3pt, dashed] (0,0) -- (2.6,1.85);
  \end{tikzpicture}}
\end{practiceq}

% ---- Q6 (point where the net force is zero, Earth-Moon)
\Needspace{12\baselineskip}
\begin{practiceq}{1}{}
The mass of Earth is about 81 times the mass of the Moon. A spacecraft travels in a straight line
from Earth to the Moon. At one point on this line, the net gravitational force on the spacecraft
(from Earth and the Moon together) is zero. Where is this point? (Distances are measured between
centres.)
\choices[1]{halfway between Earth and the Moon}
  {$\frac{1}{10}$ of the way from Earth to the Moon}
  {$\frac{9}{10}$ of the way from Earth to the Moon}
  {$\frac{81}{82}$ of the way from Earth to the Moon}
\end{practiceq}

% =====================================================================
\Needspace{32\baselineskip}
\section*{Section B: Orbits and Kepler's laws}

% ---- Q7 (forces on a satellite in a stable circular orbit)
\begin{practiceq}{1}{}\label{bq:secB}
Which diagram correctly shows \emph{all} the forces acting on a satellite in a stable circular
orbit around Earth?
\begin{center}
  \begin{tabular}{@{}L{62mm}L{62mm}@{}}
    \toprule
    $F_\mathrm{g}$ = gravitational force & $F_\mathrm{p}$ = propulsive force (engine thrust)\\
    $F_\mathrm{c}$ = centripetal force & \\
    \bottomrule
  \end{tabular}
\end{center}
\choices[2]{%
  \begin{tikzpicture}[baseline=(current bounding box.north)]
    \path (-1.8,-1.55) rectangle (1.8,2.25);
    \draw[guide] (0,0) circle (1.5);
    \shade[ball color=hscbrand!30] (0,0) circle (0.62);
    \node[font=\footnotesize] at (0,0) {Earth};
    \draw[force] (0,1.5) -- (0,0.82) node[pos=0.6, left, text=hscforce] {$F_\mathrm{g}$};
    \draw[object] (-0.15,1.42) rectangle (0.15,1.58);
  \end{tikzpicture}}
  {%
  \begin{tikzpicture}[baseline=(current bounding box.north)]
    \path (-1.8,-1.55) rectangle (1.8,2.25);
    \draw[guide] (0,0) circle (1.5);
    \shade[ball color=hscbrand!30] (0,0) circle (0.62);
    \node[font=\footnotesize] at (0,0) {Earth};
    \draw[force] (0,1.5) -- (0,0.82) node[pos=0.6, left, text=hscforce] {$F_\mathrm{g}$};
    \draw[force] (0,1.5) -- (0,2.18) node[pos=0.6, left, text=hscforce] {$F_\mathrm{c}$};
    \draw[object] (-0.15,1.42) rectangle (0.15,1.58);
  \end{tikzpicture}}
  {%
  \begin{tikzpicture}[baseline=(current bounding box.north)]
    \path (-1.8,-1.55) rectangle (1.8,2.25);
    \draw[guide] (0,0) circle (1.5);
    \shade[ball color=hscbrand!30] (0,0) circle (0.62);
    \node[font=\footnotesize] at (0,0) {Earth};
    \draw[force] (0,1.5) -- (0,0.82) node[pos=0.6, left, text=hscforce] {$F_\mathrm{g}$};
    \draw[force] (0,1.5) -- (0.8,1.5) node[above, text=hscforce] {$F_\mathrm{p}$};
    \draw[object] (-0.15,1.42) rectangle (0.15,1.58);
  \end{tikzpicture}}
  {%
  \begin{tikzpicture}[baseline=(current bounding box.north)]
    \path (-1.8,-1.55) rectangle (1.8,2.25);
    \draw[guide] (0,0) circle (1.5);
    \shade[ball color=hscbrand!30] (0,0) circle (0.62);
    \node[font=\footnotesize] at (0,0) {Earth};
    \draw[force] (-0.09,1.5) -- (-0.09,0.82) node[pos=0.6, left, text=hscforce] {$F_\mathrm{g}$};
    \draw[force] (0.09,1.5) -- (0.09,0.82) node[pos=0.6, right, text=hscforce] {$F_\mathrm{c}$};
    \draw[object] (-0.15,1.42) rectangle (0.15,1.58);
  \end{tikzpicture}}
\end{practiceq}

% ---- Q8 (geostationary orbit)
\Needspace{12\baselineskip}
\begin{practiceq}{1}{}
Which statement about a geostationary satellite is correct?
\choices[1]{It can stay above any point on Earth, for example above Sydney.}
  {It has a lower altitude than a satellite in low Earth orbit.}
  {It moves faster than a satellite in low Earth orbit.}
  {It orbits above the equator, in the same direction as Earth spins, with a period of about
   24 hours.}
\end{practiceq}

% ---- Q9 (two satellites in the same orbit, different masses)
\Needspace{10\baselineskip}
\begin{practiceq}{1}{}
Two satellites, P and Q, move in the same circular orbit around Earth. Satellite Q has three
times the mass of satellite P. Which quantity is the \emph{same} for both satellites?
\choices[2]{Kinetic energy}{Gravitational force on the satellite}
  {Gravitational potential energy}{Orbital speed}
\end{practiceq}

% ---- Q10 (orbital speed when r becomes 9r)
\Needspace{8\baselineskip}
\begin{practiceq}{1}{}
A satellite moves in a circular orbit of radius $r$ around a planet, with orbital speed $v$. It is
moved to a new circular orbit of radius $9r$. What is its new orbital speed?
\choices{$\dfrac{v}{9}$}{$\dfrac{v}{3}$}{$3v$}{$9v$}
\end{practiceq}

% ---- Q11 (period when r becomes 4r)
\Needspace{8\baselineskip}
\begin{practiceq}{1}{}
A moon orbits a planet in a circular orbit with period $T$. A second moon orbits the same planet.
Its orbital radius is 4 times as big. What is the period of the second moon?
\choices{$2T$}{$4T$}{$8T$}{$16T$}
\end{practiceq}

% ---- Q12 (which graph is a straight line)
\Needspace{10\baselineskip}
\begin{practiceq}{1}{}
For satellites in circular orbits around the same planet, $\dfrac{r^3}{T^2} = \dfrac{GM}{4\pi^2}$.
A student wants a graph that is a straight line through the origin. Which graph should the student
plot?
\choices[2]{$T$ against $r$}{$T$ against $r^3$}{$T^2$ against $r$}{$T^2$ against $r^3$}
\end{practiceq}

% ---- Q13 (Kepler's second law, equal areas)
% Ellipse a = 3, b = 2.4, e = 0.6, star at the focus (1.8, 0). The shaded sectors have equal
% areas: E from -0.5 to 0.5 rad and from pi-0.1328 to pi+0.1328 rad (numbers.py).
\Needspace{30\baselineskip}
\begin{practiceq}{1}{}
The diagram shows the elliptical orbit of a planet around a star. The two shaded areas are equal.
The planet takes 40 days to move from $P_1$ to $P_2$.
\begin{center}
  \begin{tikzpicture}[x=0.9cm, y=0.9cm]
    \fill[hscseriesA!22] (1.8,0) -- plot[domain=-0.5:0.5, samples=30, smooth]
      ({3*cos(\x r)}, {2.4*sin(\x r)}) -- cycle;
    \fill[hscseriesA!22] (1.8,0) -- plot[domain=3.0088:3.2744, samples=10, smooth]
      ({3*cos(\x r)}, {2.4*sin(\x r)}) -- cycle;
    \draw[line width=0.9pt, hscink] (0,0) ellipse (3 and 2.4);
    \draw[->, line width=0.9pt, hscink] (0.35,2.4) -- (-0.35,2.4);
    \shade[ball color=hscacc!80] (1.8,0) circle (0.2);
    \node[below=5pt, font=\small] at (1.8,0) {Star};
    \fill[hscink] (2.633,-1.151) circle (2pt) node[below right, font=\small] {$P_1$};
    \fill[hscink] (2.633,1.151) circle (2pt) node[above right, font=\small] {$P_2$};
    \fill[hscink] (-2.974,0.318) circle (2pt) node[above left, font=\small] {$P_3$};
    \fill[hscink] (-2.974,-0.318) circle (2pt) node[below left, font=\small] {$P_4$};
  \end{tikzpicture}
\end{center}
Which statement is correct?
\choices[1]{It takes 40 days to move from $P_3$ to $P_4$. The planet moves faster between $P_1$
   and $P_2$.}
  {It takes 40 days to move from $P_3$ to $P_4$. The planet moves faster between $P_3$ and $P_4$.}
  {It takes more than 40 days to move from $P_3$ to $P_4$. The planet moves faster between $P_1$
   and $P_2$.}
  {It takes less than 40 days to move from $P_3$ to $P_4$. The planet moves faster between $P_3$
   and $P_4$.}
\end{practiceq}

% ---- Q14 (Kepler's third law in AU and years)
\Needspace{9\baselineskip}
\begin{practiceq}{1}{}
A dwarf planet orbits the Sun with a period of 125 years. The average distance between Earth and
the Sun is 1 astronomical unit (AU). What is the average distance between the dwarf planet and the
Sun?
\choices{\qty{5.0}{AU}}{\qty{11}{AU}}{\qty{25}{AU}}{\qty{1400}{AU}}
\end{practiceq}

% ---- Q15 (orbital speed of the ISS)
\Needspace{8\baselineskip}
\begin{practiceq}{1}{}
The International Space Station (ISS) orbits Earth in a circular orbit at an altitude of
\qty{420}{km}. What is its orbital speed?
\choices{\qty{7.68e3}{m.s^{-1}}}{\qty{7.93e3}{m.s^{-1}}}{\qty{1.09e4}{m.s^{-1}}}
  {\qty{3.09e4}{m.s^{-1}}}
\end{practiceq}

% ---- Q16 (altitude of a geostationary satellite)
\Needspace{8\baselineskip}
\begin{practiceq}{1}{}
A geostationary satellite has an orbital period of 24.0 hours. What is its altitude above Earth's
surface?
\choices{\qty{1.80e5}{m}}{\qty{3.59e7}{m}}{\qty{4.23e7}{m}}{\qty{7.57e22}{m}}
\end{practiceq}

% ---- Q17 (mass of Jupiter from Io)
\Needspace{9\baselineskip}
\begin{practiceq}{1}{}
Io is a moon of Jupiter. Its orbit is close to a circle, with radius \qty{4.22e8}{m} and period
\qty{1.53e5}{s}. What is the mass of Jupiter?
\choices{\qty{4.50e18}{kg}}{\qty{4.81e25}{kg}}{\qty{1.90e27}{kg}}{\qty{2.91e32}{kg}}
\end{practiceq}

% ---- Q18 (T^2 against r^3 for the moons of two planets)
\Needspace{28\baselineskip}
\begin{practiceq}{1}{}
The graph shows $T^2$ against $r^3$ for the moons of two planets, X and Y.
\begin{center}
  \begin{tikzpicture}
    \draw[axisline] (0,0) -- (4.4,0) node[below left, font=\small] {$r^3$};
    \draw[axisline] (0,0) -- (0,3.0) node[above, font=\small] {$T^2$};
    \draw[hscseriesA, line width=1.3pt] (0,0) -- (2.4,2.6) node[right, text=hscink, font=\small] {Planet X};
    \draw[hscseriesB, line width=1.3pt] (0,0) -- (3.8,1.5) node[right, text=hscink, font=\small] {Planet Y};
    \foreach \x/\y in {0.6/0.65, 1.2/1.3, 1.8/1.95} {\fill[hscseriesA] (\x,\y) circle (2pt);}
    \foreach \x/\y in {0.9/0.355, 1.9/0.75, 2.85/1.125} {\fill[hscseriesB] (\x,\y) circle (2pt);}
  \end{tikzpicture}
\end{center}
Which statement is correct?
\choices[1]{Planet X has the greater mass, because its line is steeper.}
  {The two planets have the same mass, because both lines pass through the origin.}
  {The masses cannot be compared, because the moons have different masses.}
  {Planet Y has the greater mass, because its line is less steep.}
\end{practiceq}

% =====================================================================
\section*{Section C: Energy in orbits and escape velocity}

% ---- Q19 (comet in an elliptical orbit)
\Needspace{20\baselineskip}
\begin{practiceq}{1}{}\label{bq:secC}
A comet moves in an elliptical orbit around a star. $P$ is the point closest to the star. $Q$ is the
point furthest from the star.
\begin{center}
  \begin{tikzpicture}[x=0.85cm, y=0.85cm]
    \draw[line width=0.9pt, hscink] (0,0) ellipse (3.2 and 1.6);
    \draw[->, line width=0.9pt, hscink] (0.35,1.6) -- (-0.35,1.6);
    \shade[ball color=hscacc!80] (-2.77,0) circle (0.17);
    \node[right=4pt, font=\small] at (-2.77,0) {Star};
    \fill[hscink] (-3.2,0) circle (2pt) node[left, font=\small] {$P$};
    \fill[hscink] (3.2,0) circle (2pt) node[right, font=\small] {$Q$};
  \end{tikzpicture}
\end{center}
Which quantity is \emph{greater} at $Q$ than at $P$?
\choices[2]{Gravitational potential energy}{Kinetic energy}{Speed}{Total energy}
\end{practiceq}

% ---- Q20 (elliptical orbit: speed and total energy)
\Needspace{30\baselineskip}
\begin{practiceq}{1}{}
A satellite orbits Earth in an elliptical orbit that passes through positions $X$ and $Y$.
\begin{center}
  \begin{tikzpicture}[x=0.85cm, y=0.85cm]
    \draw[guide, line width=0.8pt] (0,0) ellipse (3.2 and 2.0);
    \shade[ball color=hscbrand!30] (-2.5,0) circle (0.42);
    \node[below=9pt, font=\footnotesize] at (-2.5,0) {Earth};
    \fill[hscink] (-3.2,0) circle (2pt) node[left, font=\small] {$X$};
    \fill[hscink] (3.2,0) circle (2pt) node[right, font=\small] {$Y$};
  \end{tikzpicture}
\end{center}
Which row of the table is correct?
\begin{center}
  \renewcommand\arraystretch{1.45}
  \begin{tabular}{@{}C{10mm}C{38mm}C{55mm}@{}}
    \toprule
    & \colhead{Greater speed} & \colhead{Greater total energy (kinetic + potential)}\\
    \midrule
    \bkchoice{A} & $X$ & $X$\\
    \bkchoice{B} & $X$ & the same at $X$ and $Y$\\
    \bkchoice{C} & $Y$ & $Y$\\
    \bkchoice{D} & $Y$ & the same at $X$ and $Y$\\
    \bottomrule
  \end{tabular}
\end{center}
\end{practiceq}

% ---- Q21 (LEO and GEO: potential energy and speed)
\Needspace{18\baselineskip}
\begin{practiceq}{1}{}
Two satellites have the same mass. One is in low Earth orbit (LEO). The other is in geostationary
orbit (GEO). Which row of the table is correct?
\begin{center}
  \renewcommand\arraystretch{1.45}
  \begin{tabular}{@{}C{10mm}C{60mm}C{40mm}@{}}
    \toprule
    & \colhead{Greater gravitational potential energy} & \colhead{Greater orbital speed}\\
    \midrule
    \bkchoice{A} & LEO & LEO\\
    \bkchoice{B} & LEO & GEO\\
    \bkchoice{C} & GEO & LEO\\
    \bkchoice{D} & GEO & GEO\\
    \bottomrule
  \end{tabular}
\end{center}
\end{practiceq}

% ---- Q22 (moving to a higher orbit)
\Needspace{11\baselineskip}
\begin{practiceq}{1}{}
A satellite is moved from a low circular orbit to a higher circular orbit around Earth. Which
statement is correct?
\choices[1]{Its kinetic energy increases.}{Its gravitational potential energy decreases.}
  {Its orbital period decreases.}{Its total energy increases.}
\end{practiceq}

% ---- Q23 (potential energy of a satellite)
\Needspace{8\baselineskip}
\begin{practiceq}{1}{}
A satellite of mass \qty{1.20e3}{kg} orbits Earth at an altitude of \qty{3.50e3}{km}. What is its
gravitational potential energy?
\choices{\qty{-2.43e10}{J}}{\qty{-4.87e10}{J}}{\qty{-7.54e10}{J}}{\qty{-1.37e11}{J}}
\end{practiceq}

% ---- Q24 (change in U read from a graph; m = 200 kg)
\Needspace{40\baselineskip}
\begin{practiceq}{1}{}
The graph shows how the gravitational potential energy $U$ of a \qty{200}{kg} satellite depends on
its altitude above Earth's surface.
\begin{center}
  \begin{tikzpicture}
    \begin{axis}[hsc, width=118mm, height=92mm,
        axis x line*=top, axis y line*=left,
        xmin=0, xmax=20000, ymin=-13, ymax=0,
        xtick={0,4000,...,20000}, ytick={-13,-12,...,0},
        minor x tick num=3, minor y tick num=1,
        grid=both, minor grid style={hscrule!45, line width=0.2pt},
        xticklabel style={/pgf/number format/1000 sep={\,}, font=\footnotesize},
        xlabel={Altitude (km)}, ylabel={$U$ ($\times 10^{9}$ J)},
        xlabel style={yshift=2pt}]
      \addplot[hscseriesA, line width=1.4pt, domain=0:20000, samples=120, smooth]
        {-6.67e-11*6.0e24*200/((6371+x)*1000)/1e9};
    \end{axis}
  \end{tikzpicture}
\end{center}
The satellite is moved from an altitude of \qty{2000}{km} to an altitude of \qty{12000}{km}. What is
the change in its gravitational potential energy?
\choices{\qty{-1.4e10}{J}}{\qty{-5.2e9}{J}}{\qty[retain-explicit-plus]{+5.2e9}{J}}{\qty[retain-explicit-plus]{+1.4e10}{J}}
\end{practiceq}

% ---- Q25 (force falls to one-ninth: what happens to U)
\Needspace{12\baselineskip}
\begin{practiceq}{1}{}
A satellite is moved to a new circular orbit. In the new orbit, the gravitational force on the
satellite is one-ninth of the force in its first orbit. What is its gravitational potential
energy in the new orbit? (Both values of potential energy are negative.)
\choices[1]{one-third of its value in the first orbit}{one-ninth of its value in the first orbit}
  {three times its value in the first orbit}{nine times its value in the first orbit}
\end{practiceq}

% ---- Q26 (escape velocity: distance matters, mass does not)
\Needspace{10\baselineskip}
\begin{practiceq}{1}{}
A space probe of mass $m$ is at a distance $8r$ from the centre of Earth. The escape velocity for
this probe is $v$. A second probe, of mass $2m$, is at a distance $2r$ from the centre of Earth.
What is the escape velocity for the second probe?
\choices{$\dfrac{v}{2}$}{$\sqrt{2}\,v$}{$2v$}{$4v$}
\end{practiceq}

% ---- Q27 (rank escape velocities from a data table)
\Needspace{20\baselineskip}
\begin{practiceq}{1}{}
The table gives data for three planets, X, Y and Z.
\begin{center}
  \begin{tabular}{@{}C{18mm}C{30mm}C{30mm}C{42mm}@{}}
    \toprule
    \colhead{Planet} & \colhead{Mass (kg)} & \colhead{Diameter (km)} &
      \colhead{Distance from its star ($\times 10^{6}$ km)}\\
    \midrule
    X & \num{2.0e24} & \num{8000} & 80\\
    Y & \num{6.0e24} & \num{20000} & 150\\
    Z & \num{1.0e24} & \num{3000} & 300\\
    \bottomrule
  \end{tabular}
\end{center}
Which list gives the escape velocities from the surfaces of the planets, from largest to smallest?
\choices{Z, Y, X}{Y, X, Z}{Z, X, Y}{X, Y, Z}
\end{practiceq}

% ---- Q28 (escape velocity from Mercury)
\Needspace{8\baselineskip}
\begin{practiceq}{1}{}
Mercury has a mass of \qty{3.30e23}{kg} and a radius of \qty{2.44e6}{m}. What is the escape velocity
from the surface of Mercury?
\choices{\qty{3.00e3}{m.s^{-1}}}{\qty{4.25e3}{m.s^{-1}}}{\qty{1.34e5}{m.s^{-1}}}
  {\qty{1.80e7}{m.s^{-1}}}
\end{practiceq}

% ---- Q29 (U of a mass moved along the line between two identical planets)
% Correct curve: U = -k (1/x + 1/(3 - x)) for x from 0.25 to 2.75, with k = 0.45.
\Needspace{36\baselineskip}
\begin{practiceq}{1}{}
Two identical planets, P and Q, are a fixed distance apart. A small mass is moved at a steady speed
along the straight line from point $S$ (near P) to point $F$ (near Q). Ignore all other masses.
\begin{center}
  \begin{tikzpicture}
    \shade[ball color=hscbrand!30] (0,0) circle (0.55);
    \node[font=\small] at (0,0) {P};
    \shade[ball color=hscbrand!30] (6,0) circle (0.55);
    \node[font=\small] at (6,0) {Q};
    \draw[displacement] (0.95,0) -- (5.05,0);
    \fill[object] (0.95,0) circle (3pt) node[above=3pt, font=\small, text=hscink] {$S$};
    \fill[hscink] (5.05,0) circle (1.5pt) node[above=3pt, font=\small] {$F$};
  \end{tikzpicture}
\end{center}
Which graph best shows the gravitational potential energy $U$ of the small mass during this
movement?
\choices[2]{%
  \begin{tikzpicture}[baseline=(current bounding box.north)]
    \draw[axisline] (0,0) -- (3.2,0) node[above left, font=\small] {Time};
    \draw[axisline] (0,0) -- (0,-2.3) node[below, font=\small] {$U$};
    \node[left, font=\footnotesize, text=hscmuted] at (0,0) {0};
    \draw[hscseriesA, line width=1.3pt] (0.25,-1.3) -- (2.75,-1.3);
  \end{tikzpicture}}
  {%
  \begin{tikzpicture}[baseline=(current bounding box.north)]
    \draw[axisline] (0,0) -- (3.2,0) node[above left, font=\small] {Time};
    \draw[axisline] (0,0) -- (0,-2.3) node[below, font=\small] {$U$};
    \node[left, font=\footnotesize, text=hscmuted] at (0,0) {0};
    \draw[hscseriesA, line width=1.3pt, domain=0.25:2.75, samples=60, smooth]
      plot (\x, {-2.56 + 0.45*(1/\x + 1/(3-\x))});
  \end{tikzpicture}}
  {%
  \begin{tikzpicture}[baseline=(current bounding box.north)]
    \draw[axisline] (0,0) -- (3.2,0) node[above left, font=\small] {Time};
    \draw[axisline] (0,0) -- (0,-2.3) node[below, font=\small] {$U$};
    \node[left, font=\footnotesize, text=hscmuted] at (0,0) {0};
    \draw[hscseriesA, line width=1.3pt] (0.25,-1.96) -- (2.75,-0.6);
  \end{tikzpicture}}
  {%
  \begin{tikzpicture}[baseline=(current bounding box.north)]
    \draw[axisline] (0,0) -- (3.2,0) node[above left, font=\small] {Time};
    \draw[axisline] (0,0) -- (0,-2.3) node[below, font=\small] {$U$};
    \node[left, font=\footnotesize, text=hscmuted] at (0,0) {0};
    \draw[hscseriesA, line width=1.3pt, domain=0.25:2.75, samples=60, smooth]
      plot (\x, {-0.45*(1/\x + 1/(3-\x))});
  \end{tikzpicture}}
\end{practiceq}

% ---- Q30 (minimum energy to move to a higher orbit)
\Needspace{12\baselineskip}
\begin{practiceq}{1}{}
An \qty{800}{kg} satellite is moved from a circular orbit of radius \qty{7.0e6}{m} to a circular
orbit of radius \qty{2.1e7}{m} around Earth. The total energy of a satellite in a circular orbit is
$E = -\dfrac{GMm}{2r}$. What is the minimum energy that must be supplied to the satellite?
\choices{\qty{7.6e9}{J}}{\qty{1.5e10}{J}}{\qty{3.0e10}{J}}{\qty{4.6e10}{J}}
\end{practiceq}
